% Part: first-order-logic
% Chapter: syntax-and-semantics
% Section: formation-sequences

\documentclass[../../../include/open-logic-section]{subfiles}

\begin{document}

\olfileid{pl}{syn}{fseq}

\olsection{Vormingsrijen}

Met een inductieve definitie kunnen we !!{formula}s op een elegante en
doelmatige manier definiëren. Met de bijbehorende inductiemethode
kunnen we eigenschappen van !!{formula}s bewijzen. Soms is het ook
nuttig om van onderaf, stap voor stap, te bekijken hoe we !!{formula}s
opbouwen. Daarvoor gebruiken we hier een \emph{vormingsrij}.

\begin{defn}[Vormingsrijen voor formules]
\ollabel{defn:fseq-frm}
Neem een eindige rij $\tuple{!A_0,\dotsc,!A_n}$ van tekenreeksen uit de
taal~$\Lang L_0$. Deze rij is een \emph{vormingsrij} voor $!A$ als
$!A \ident !A_n$ en als voor elke $i \leq n$ een van twee dingen
geldt. Ofwel is $!A_i$ een atomaire formule, ofwel bestaan er
$j,k < i$ waarvoor een van de volgende gevallen geldt:
\begin{enumerate}
    \tagitem{prvNot}{$!A_i \ident \lnot !A_j$.}{}%
    \tagitem{prvAnd}{$!A_i \ident (!A_j \land !A_k)$.}{}%
    \tagitem{prvOr}{$!A_i \ident (!A_j \lor !A_k)$.}{}%
    \tagitem{prvIf}{$!A_i \ident (!A_j \lif !A_k)$.}{}%
    \tagitem{prvIff}{$!A_i \ident (!A_j \liff !A_k)$.}{}%
\end{enumerate}
\end{defn}

\begin{ex}
\[
\tuple{
    \Obj p_0,
    \Obj p_1,
    (\Obj p_1 \land \Obj p_0),
    \lnot (\Obj p_1 \land \Obj p_0)
}
\]
is een vormingsrij van
$\lnot (\Obj p_1 \land \Obj p_0)$. De volgende rij is dat ook:
\[
\tuple{
    \Obj p_0,
    \Obj p_1,
    \Obj p_0,
    (\Obj p_1 \land \Obj p_0),
    (\Obj p_0 \lif \Obj p_1),
    \lnot (\Obj p_1 \land \Obj p_0)
}.
\]
%
Het tweede voorbeeld laat zien dat een vormingsrij ``rommel'' mag
bevatten. Er mogen formules in staan die overbodig zijn of niet worden
gebruikt om de laatste formule op te bouwen.
\end{ex}

\begin{prop}
\ollabel{prop:formed}
Elke !!{formula}~$!A$ in~$\Frm[L_0]$ heeft een vormingsrij.
\end{prop}

\begin{proof}
Neem eerst aan dat $!A$ atomair is. Dan is de rij~$\tuple{!A}$ een
vormingsrij voor~$!A$.
%
Neem nu aan dat $!B$ en~$!C$ vormingsrijen hebben. Noem die rijen
$\tuple{!B_0,\dotsc,!B_n}$ en $\tuple{!C_0,\dotsc,!C_m}$.
%
\begin{enumerate}
  \tagitem{prvNot}{Als $!A \ident \lnot !B$, dan is
      $\tuple{!B_0,\dotsc,!B_n,\lnot !B_n}$
      een vormingsrij voor~$!A$.}{}
  \tagitem{prvAnd}{Als $!A \ident (!B \land !C)$, dan is
      $\tuple{!B_0,\dotsc,!B_n,!C_0,\dotsc,!C_m,(!B_n \land !C_m)}$
      een vormingsrij voor~$!A$.}{}
  \tagitem{prvOr}{Als $!A \ident (!B \lor !C)$, dan is
      $\tuple{!B_0,\dotsc,!B_n,!C_0,\dotsc,!C_m,(!B_n \lor !C_m)}$
      een vormingsrij voor~$!A$.}{}
  \tagitem{prvIf}{Als $!A \ident (!B \lif !C)$, dan is
      $\tuple{!B_0,\dotsc,!B_n,!C_0,\dotsc,!C_m,(!B_n \lif !C_m)}$
      een vormingsrij voor~$!A$.}{}
  \tagitem{prvIff}{Als $!A \ident (!B \liff !C)$, dan is
      $\tuple{!B_0,\dotsc,!B_n,!C_0,\dotsc,!C_m,(!B_n \liff !C_m)}$
      een vormingsrij voor~$!A$.}{}
\end{enumerate}
Het inductieprincipe voor !!{formula}s laat nu zien dat elke
!!{formula} een vormingsrij heeft.
\end{proof}

Ook het omgekeerde geldt. Dat is belangrijk, want daarmee weten we dat
onze twee manieren om formules te definiëren hetzelfde resultaat
geven. Bovendien kunnen we daardoor stellingen over formules bewijzen
met gewone inductie naar de lengte van vormingsrijen.

\begin{lem}
\ollabel{lem:fseq-init}
Neem aan dat $\tuple{!A_0,\dotsc,!A_n}$ een vormingsrij voor~$!A_n$ is
en dat $k \leq n$. Dan is $\tuple{!A_0,\dotsc,!A_k}$ een vormingsrij
voor~$!A_k$.
\end{lem}

\begin{proof}
Opgave.
\end{proof}

\begin{thm}
\ollabel{thm:fseq-frm-equiv}
$\Frm[L_0]$ is precies de verzameling van alle tekenreeksen in de
taal~$\Lang L_0$ die een vormingsrij hebben.
\end{thm}

\begin{proof}
  Noem $F$ de verzameling van alle tekenreeksen in de taal~$\Lang L_0$
  die een vormingsrij hebben. Uit
\olref[pl][syn][fseq]{prop:formed} weten we al dat
$\Frm[L_0] \subseteq F$. Nu bewijzen we de andere richting.

  Neem aan dat $!A$ de vormingsrij $\tuple{!A_0,\dotsc,!A_n}$ heeft. We
  bewijzen dat $!A \in \Frm[L_0]$ met sterke inductie naar~$n$. De
inductiehypothese is deze: elke tekenreeks met een vormingsrij waarvan
de laatste index $m < n$ is, behoort tot~$\Frm[L_0]$. Uit de definitie
van een vormingsrij volgt nu een tweedeling. Ofwel is $!A_n$ atomair,
ofwel bestaan er $j,k < n$ waarvoor een van de volgende gevallen
geldt:
\begin{enumerate}
    \tagitem{prvNot}{$!A_n \ident \lnot !A_j$.}{}%
    \tagitem{prvAnd}{$!A_n \ident (!A_j \land !A_k)$.}{}%
    \tagitem{prvOr}{$!A_n \ident (!A_j \lor !A_k)$.}{}%
    \tagitem{prvIf}{$!A_n \ident (!A_j \lif !A_k)$.}{}%
    \tagitem{prvIff}{$!A_n \ident (!A_j \liff !A_k)$.}{}%
\end{enumerate}
We bekijken de gevallen afzonderlijk. Als $!A_n$ atomair is, dan is
$!A_n \in \Frm[L_0]$. Neem nu het geval $!A \ident
(!A_j \land !A_k)$. Volgens \olref[pl][syn][fseq]{lem:fseq-init} zijn
$\tuple{!A_0,\dotsc,!A_j}$ en $\tuple{!A_0,\dotsc,!A_k}$ vormingsrijen
voor $!A_j$ en~$!A_k$. Beide rijen zijn echte beginstukken van de
vormingsrij voor~$!A$. Hun laatste indices zijn dus kleiner dan~$n$.
Volgens de inductiehypothese behoren $!A_j$ en~$!A_k$ dus
  tot~$\Frm[L_0]$. Volgens de definitie van !!a{formula} geldt dat daarom
  ook voor $(!A_j \land !A_k)$. De andere gevallen gaan op
dezelfde manier.
\end{proof}
\end{document}
